% ============================================================
\subsubsection{Column Operations and Vectorization}
% ============================================================

pandas operations run array-at-a-time. A column expression applies to the whole Series in one pass over the underlying array, which runs far faster than a Python loop and removes the need for a loop in most transformations. This vectorized style is the default, and the row-wise \texttt{apply} is the fallback for the rare case that resists it.

\paragraph{Building columns.} Arithmetic and comparisons between columns produce new Series aligned on the index, assigned back by name. \texttt{assign} does the same inside a chain, returning a copy with columns added or overwritten. Its callable form receives the frame, which lets a new column read from the frame as it is being built.

\begin{lstlisting}[style=python, caption={Direct assignment and the chainable assign}]
df['ratio'] = df['rev'] / df['units']

df = df.assign(
    ratio=lambda d: d['rev'] / d['units'],
    is_big=lambda d: d['rev'] > 1000)
\end{lstlisting}

\begin{keybox}[Prefer vectorized operations over apply]
Arithmetic, comparisons, numpy ufuncs, and the \texttt{.str} and \texttt{.dt} accessors are all vectorized and act on the whole column at compiled speed. A Python-level \texttt{apply} runs the function once per element or row and is slower by one to two orders of magnitude. Nearly every transformation has a vectorized form, and finding it repays the effort at any nontrivial size.
\end{keybox}

\vspace{1ex}

\paragraph{Conditional columns.} A two-way choice is \texttt{np.where(cond, a, b)}, taking \texttt{a} where the condition holds and \texttt{b} elsewhere. A multi-way choice is \texttt{np.select(conditions, choices, default)}, which takes the first matching condition in order. \texttt{Series.where(cond, other)} keeps the value where the condition holds and substitutes \texttt{other} elsewhere, and \texttt{mask} is its complement. Each of these replaces a conditional \texttt{apply} with a single vectorized call.

\begin{lstlisting}[style=python, caption={Two-way and multi-way conditionals without apply}]
df['flag'] = np.where(df['temp'] > 97, 'hot', 'mild')

conditions = [df['score'] >= 90, df['score'] >= 80]
grades     = ['A', 'B']
df['grade'] = np.select(conditions, grades, default='C')
\end{lstlisting}

\vspace{-1ex}

\paragraph{The str and dt accessors.} A Series of string dtype exposes vectorized string methods under \texttt{.str}, and a datetime Series exposes calendar and clock fields under \texttt{.dt}. Both return aligned Series, so their results drop straight into masks and new columns. \texttt{pd.to\_datetime} parses text or epoch values into \texttt{datetime64[ns]}, the prerequisite for \texttt{.dt}.

\begin{lstlisting}[style=python, caption={Vectorized string and datetime work}]
df['user'].str.lower().str.startswith('a')      # chainable, vectorized
df['email'].str.contains('@company', na=False)  # na=False maps missing to False

df['ts']      = pd.to_datetime(df['ts'])
df['year']    = df['ts'].dt.year
df['weekday'] = df['ts'].dt.dayofweek           # Monday = 0
\end{lstlisting}

\vspace{-1ex}

The workhorses are \texttt{.str.contains}, \texttt{.str.extract} for a regex capture, \texttt{.str.split}, \texttt{.str.len}, and positional slicing through \texttt{.str[:3]}, alongside \texttt{.dt.year}, \texttt{.dt.month}, \texttt{.dt.date}, and \texttt{.dt.floor('D')}. The \texttt{na} argument on the string predicates fixes how a missing value scores in the resulting mask.

\paragraph{Element mapping.} \texttt{map} on a Series applies a dict or function element by element, which suits recoding a small fixed set of values. \texttt{replace} substitutes chosen values throughout a column, \texttt{astype} casts to a target dtype, and \texttt{round} sets the number of decimals on a numeric column.

\paragraph{Binning into categories.} \texttt{pd.cut} slices a numeric column at explicit edges into labeled bands, the vectorized form of a \texttt{CASE} ladder over ranges. \texttt{pd.qcut} instead cuts at quantiles, producing bins of roughly equal count, which is the \texttt{NTILE} operation. Both return a \texttt{category} column that feeds \texttt{value\_counts} or a \texttt{groupby} directly.

\newpage

\begin{lstlisting}[style=python, caption={Fixed-edge bins and equal-frequency bins}]
df['band'] = pd.cut(df['salary'],
                    bins=[0, 50_000, 80_000, np.inf],
                    labels=['<50k', '50-80k', '80k+'])

df['quartile'] = pd.qcut(df['salary'], 4, labels=False)  # 0..3 by quantile
\end{lstlisting}

\vspace{1ex}

\begin{keybox}[apply and the axis argument]
\texttt{Series.apply(f)} calls \texttt{f} once per element, and \texttt{DataFrame.apply(f, axis=1)} calls it once per row, handing each row over as a Series. Both run at Python speed and are the last resort, kept for logic with no vectorized form. The \texttt{axis} argument is a recurring source of confusion: \texttt{axis=0} runs down the index and operates on each column, while \texttt{axis=1} runs across the columns and operates on each row. A row-wise \texttt{apply} spanning several columns is the one case where \texttt{axis=1} earns its cost.
\end{keybox}
